\booksection{Kapitel 11 Anhang: Professionelle Bildtechnik}

\chapterintro{Um wirklich die Geheimnisse der Smartphone-Fotografie und -Videografie zu verstehen, sollten Sie einige professionelle Bildtechnik-Kenntnisse kennenlernen. Dieses Wissen hilft Ihnen nicht nur, die Smartphone-Kamera besser zu nutzen, sondern verleiht Ihnen auch bei der Geräteauswahl mehr Urteilskraft.}

\booksubsection{Auflösung, Bildrate und Bitrate: Die drei Elemente der Videoqualität}

\subsubsection*{Auflösung (Resolution):}

Die Auflösung bestimmt die Schärfe des Bildes und wird in Pixeln angegeben.

\begin{itemize}
  \item \textbf{720P (HD)}: 1280×720 Pixel, Einsteiger-Schärfe.
  \item \textbf{1080P (Full HD)}: 1920×1080 Pixel, derzeit die gängigste Auflösung.
  \item \textbf{2K (QHD)}: 2560×1440 Pixel, schärfer als 1080P.
  \item \textbf{4K (UHD)}: 3840×2160 Pixel, Ultra HD, detailreich.
  \item \textbf{8K}: 7680×4320 Pixel, Zukunftsstandard, derzeit begrenzte Geräteunterstützung.
\end{itemize}

Je höher die Auflösung, desto schärfer das Bild, aber die Datei wird größer und die Anforderungen an Speicher und Rechenleistung steigen.

\subsubsection*{Bildrate (Frame Rate):}

Die Bildrate bestimmt die Flüssigkeit der Darstellung und wird in Bildern pro Sekunde (fps) angegeben.

\begin{itemize}
  \item \textbf{24fps}: Standard-Bildrate des Kinos, mit eigenem „Film-Look".
  \item \textbf{30fps}: Für TV und Webvideos üblich, flüssig und natürlich.
  \item \textbf{60fps}: Hohe Bildrate, ideal für Sport und Gaming-Aufnahmen, sehr flüssig.
  \item \textbf{120fps/240fps}: Sehr hohe Bildrate, für Zeitlupen-Videos.
\end{itemize}

Je höher die Bildrate, desto flüssiger das Bild und desto besser der Zeitlupeneffekt, aber die Datei wird ebenfalls größer.

\subsubsection*{Bitrate (Bitrate):}

Die Bitrate bestimmt die Datendichte des Videos und wird in Datenmenge pro Sekunde (Mbps) angegeben.

\begin{itemize}
  \item Höhere Bitrate bedeutet bessere Qualität und mehr Details, aber auch größere Dateien.
  \item Zu niedrige Bitrate führt zu Blockbildung, Farbverfälschungen und anderen Kompressionsartefakten.
  \item Professionelle Videoaufnahmen erfordern meist eine hohe Bitrate, um mehr Spielraum für die Nachbearbeitung zu behalten.
\end{itemize}

\booksubsection{Video-Codecs: H.264, H.265 und ProRes}

\subsubsection*{H.264 (AVC):}

\begin{itemize}
  \item \textbf{Merkmale}: Der am weitesten verbreitete Codec, extrem gute Kompatibilität.
  \item \textbf{Vorteile}: Fast alle Geräte können ihn abspielen, Kompressionseffizienz moderat.
  \item \textbf{Nachteile}: Im Vergleich zu neueren Formaten geringere Kompression, größere Dateien.
  \item \textbf{Anwendungsfall}: Teilen von Webvideos, Alltagsaufnahmen normaler Nutzer.
\end{itemize}

\subsubsection*{H.265 (HEVC):}

\begin{itemize}
  \item \textbf{Merkmale}: Die Weiterentwicklung von H.264, etwa 50 \% höhere Kompressionseffizienz.
  \item \textbf{Vorteile}: Bei gleicher Qualität kleinere Dateien, unterstützt höhere Auflösungen und Bildraten.
  \item \textbf{Nachteile}: Codierung und Decodierung erfordern stärkere Prozessoren, einige ältere Geräte sind nicht kompatibel.
  \item \textbf{Anwendungsfall}: Für 4K-Videoaufnahmen und Szenarien mit begrenztem Speicherplatz.
\end{itemize}

\subsubsection*{Apple ProRes:}

\begin{itemize}
  \item \textbf{Merkmale}: Der von Apple entwickelte professionelle Video-Codec, nahezu verlustfreie Kompression.
  \item \textbf{Vorteile}: Extrem hohe Bildqualität, großer Spielraum für die Nachbearbeitung, farbgenau.
  \item \textbf{Nachteile}: Dateien extrem groß, erfordern viel Speicherplatz und professionelle Geräte zur Bearbeitung.
  \item \textbf{Anwendungsfall}: Professionelle Videoproduktion, filmreife Aufnahmen.
\end{itemize}

\subsubsection*{ProRes RAW:}

\begin{itemize}
  \item \textbf{Merkmale}: Kombiniert die Vorteile von ProRes und RAW.
  \item \textbf{Was ist RAW?}: Unverarbeitete Original-Sensordaten, bewahren die maximale Informationsmenge.
  \item \textbf{Vorteile}: Enormer Spielraum bei der Nachbearbeitung – Belichtung, Weißabgleich und andere Parameter lassen sich ändern, ohne Qualität zu verlieren.
  \item \textbf{Nachteile}: Dateien extrem groß, erfordert professionelle Nachbearbeitungs-Software.
  \item \textbf{Anwendungsfall}: High-End-Profiproduktion. Ab der iPhone 15 Pro-Serie wird Apple ProRes-Video (bis 4K 60fps) unterstuetzt; das iPhone 17 Pro zusaetzlich mit ProRes RAW direkt im Geraet.
\end{itemize}

\booksubsection{HDR: High-Dynamic-Range-Technologie}

HDR (High Dynamic Range) kann gleichzeitig Details in den hellsten und dunkelsten Bildbereichen erhalten.

\begin{itemize}
  \item \textbf{HDR10}: Der grundlegendste HDR-Standard, gute Kompatibilität.
  \item \textbf{HDR10+}: Von Samsung vorangetriebener Standard, unterstützt dynamische Metadaten.
  \item \textbf{Dolby Vision}: Der von Dolby entwickelte Top-HDR-Standard, beste Qualität, erfordert jedoch Lizenzierung.
  \item \textbf{HLG}: Rundfunk-HDR-Standard, kompatibel mit gewöhnlichen Displays.
\end{itemize}

Bei HDR-Videoaufnahmen wird ein HDR-fähiger Bildschirm benötigt, um die volle Wirkung zu sehen.

\booksubsection{Farbraum und Farbtiefe}

\subsubsection*{Farbraum:}

\begin{itemize}
  \item \textbf{sRGB}: Standard-Farbraum, beste Kompatibilität.
  \item \textbf{Display P3}: Ein von Apple geförderter Breitfarbraum-Standard, reichere Farben.
  \item \textbf{Rec.2020}: Standard-Farbraum für HDR-Videos, breitester Farbumfang.
\end{itemize}

\subsubsection*{Farbtiefe:}

\begin{itemize}
  \item \textbf{8-bit}: 16,7 Millionen Farben, Standard für normale Videos.
  \item \textbf{10-bit}: 1,07 Milliarden Farben, Standard für professionelle Videos, sanftere Farbübergänge.
  \item \textbf{12-bit}: 68 Milliarden Farben, hochwertiger Profistandard.
\end{itemize}

\begin{tcbinfo}{Praktische Empfehlungen}
\begin{itemize}
  \item Alltag: 1080P/4K + 30fps + H.265-Codierung, ausgewogen zwischen Qualität und Dateigröße.
  \item Sport: 1080P/4K + 60fps, für flüssige Darstellung.
  \item Professionelle Produktion: 4K + ProRes (falls unterstützt), für maximalen Spielraum bei der Nachbearbeitung.
  \item Speicher-Tipp: Hochbitrate-Videos belegen viel Platz – regelmäßig auf Computer oder Cloud sichern.
\end{itemize}

\end{tcbinfo}

\begin{tcbpractice}{Entdecken Sie: Ihre Smartphone-Videoeinstellungen}
\begin{enumerate}
  \item Öffnen Sie die Kamera-App und wechseln Sie in den Videomodus.
  \item Prüfen Sie die Optionen für Auflösung und Bildrate und probieren Sie verschiedene Kombinationen.
  \item Falls unterstützt, prüfen Sie, ob HDR, ProRes und weitere Profi-Optionen vorhanden sind.
  \item Nehmen Sie einige Videos mit unterschiedlichen Einstellungen auf und vergleichen Sie Qualität und Dateigröße.
\end{enumerate}

\end{tcbpractice}

\booksubsection{V. RAW-Format im Vergleich zu HEIF/JPEG}

\subsubsection*{Was ist das RAW-Format:}

RAW ist das vom Kamerasensor erfasste „Rohdatenformat“ – unverdichtete, unbearbeitete Ursprungsinformationen. Wie frische Zutaten, die erst „zubereitet” werden muessen. Smartphones speichern standardmaessig fertig „zugerichtete“ JPEG- oder HEIF-Dateien; RAW muss explizit aktiviert werden.

\subsubsection*{RAW vs. HEIF vs. JPEG im Vergleich:}

\begin{itemize}
  \item \textbf{JPEG}: Aeltestes Format, beste Kompatibilitaet (jedes Geraet kann es oeffnen). 8-bit-Farbtiefe, verlustbehaftete Komprimierung, kleinste Dateien (2-5 MB pro Foto). Begrenzter Spielraum fuer Nachbearbeitung; starkes Aufhellen dunkler Partien erzeugt Farbblockartefakte.
  \item \textbf{HEIF/HEIC}: Das von Apple propagierte neuere Format; bei gleicher Qualitaet nur halb so gross wie JPEG (1-3 MB pro Foto). Unterstuetzt 10-bit-Farbtiefe und transparente Kanuele. Nachteil: Aeltere Geraete und Windows benoetigen Plugins zum Oeffnen.
  \item \textbf{RAW}: Unkomprimierte Rohdaten, groesste Dateien (20-100 MB pro Foto). Farbtiefe meist 12-16 bit, maximaler Spielraum fuer Nachbearbeitung. Nachteil: Benoetigt professionelle Software (Lightroom, Snapseed etc.) zum Oeffnen und Bearbeiten; nicht direkt teilbar.
\end{itemize}

\subsubsection*{Warum Profifotografen RAW aufnehmen:}

\begin{itemize}
  \item \textbf{Weissabgleich nachtraeglich verlustfrei anpassbar}: Wenn bei JPEG der Weissabgleich falsch eingestellt war (z. B. gelblich unter Gluehlampenlicht), kostet die Korrektur Qualitaet; bei RAW laesst sich die Farbtemperatur neu setzen – komplett verlustfrei.
  \item \textbf{Grosse Belichtungstoleranz}: Bei 2 Blenden Unterbelichtung ist das Dunkel bei JPEG praktisch schwarz; Aufhellen erzeugt nur Rauschen. RAW kann aus 3-4 Blenden Unterbeligung noch nutzbare Details „retten". Ueberbelichtung entsprechend.
  \item \textbf{Mehr Details in Schatten und Lichtern}Der Dynamikumfang von RAW ist wesentlich groesser als bei JPEG – mehr Spielraum, bis Lichter ausbrennen oder Schatten zulaufen.
  \item \textbf{Schaerfung, Rauschunterdrueckung und Objektivkorrektur vollstaendig nachtraeglich steuerbar}: Die kamerainterne Bearbeitung ist „nur einmal machbar"; RAW erlaubt wiederholte Anpassungen, immer ausgehend von den Originaldaten.
\end{itemize}

\subsubsection*{So aktivierst du RAW auf dem Smartphone:}

\begin{itemize}
  \item \textbf{iPhone ProRAW}: Nur ab iPhone 12 Pro und hoeher (Pro-Modelle). Unter „Einstellungen Kamera Formate" „Apple ProRAW" aktivieren. Danach erscheint oben in der Kamera-Oberflaeche eine „RAW"-Taste – antippen zum Aktivieren. ca. 25 MB pro Foto; nur bei Bedarf einschalten.
  \item \textbf{Android Motion Raw / Ultra RAW}: Unterschiedliche Bezeichnungen je Hersteller. Bei Xiaomi/OPPO/vivo-Flaggschiffen findet sich der RAW-Schalter im „Profi-Modus" oder unter „Mehr"; Huawei im „Profi-Modus"; Samsung Galaxy S-Serie im „Profi-Modus". Dateigroesse meist 20-50 MB.
  \item \textbf{Allgemeine Empfehlung}: Fuer Alltagsaufnahmen muss RAW nicht aktiviert werden – die Dateien sind zu gross; nur bei wichtigen Motiven (Landschaft, Portraet, Nachbearbeitung erforderlich) einschalten.
\end{itemize}

\booksubsection{VI. Grundlagen der Computational Photography}

Die moderne Smartphone-Fotografie ist unzertrennlich mit „Computational Photography“ verbunden – also der Erzeugung von Fotos durch Algorithmen, nicht nur durch Linse und Sensor. Das Verstaendnis dieser Kernprinzipien erleichtert es dir zu begreifen, was das Smartphone in dem Moment des Ausloesers tut.

\subsubsection*{Multi-Frame-Compositing:}

Wenn du einmal abdrueckst, macht das Smartphone in Wirklichkeit mehrere Aufnahmen. Der Algorithmus richtet die Frames aus, gewichtet und verschmilzt sie zu einem hochwertigen Endbild.

\begin{itemize}
  \item \textbf{Nachtmodus}: Innerhalb von 1-3 Sekunden werden mehrere unterbelichtete Aufnahmen gemacht, ausgerichtet und ueberlagert – Rauschen hebt sich auf, Helligkeit summiert sich. So entstehen auch bei schwachem Licht saubere Bilder.
  \item \textbf{HDR (High Dynamic Range)}: Es werden unter-, normal- und ueberbelichtete Aufnahmen gemacht (drei oder mehr) und zu einem Bild verschmolzen, in dem sowohl Licht- als auch Schattenpartien Details zeigen. Besonders effektiv bei Gegenlicht.
  \item \textbf{Deep Fusion (Apple) / Photonic Engine}: Fuer Mittel- bis Schwachlicht-Szenen; analysiert mehrere Frames pixelweise und verschmilzt sie, mit Schwerpunkt auf Textur- und Detailverbesserung – Strickpullover und Haare werden deutlich schaerfer.
\end{itemize}

\subsubsection*{Semantic Segmentation:}

Die KI verarbeitet nicht das gesamte Bild gleich, sondern „versteht“ zunaechst, was im Bild ist – Himmel, Gesichter, Wiese, Gebaeude, Kleidung – und wendet dann auf jeden Bereich unterschiedliche Verarbeitungsstrategien an.

\begin{itemize}
  \item \textbf{Himmel}: Blaetoe separat verstaerken, Helligkeit absenken, um Ueberbelichtung zu vermeiden.
  \item \textbf{Gesicht}: Hautton-Bereich erkennen, separat aufhellen und verjuengen, ohne Gesichtsmerkmale und Haardetails zu veraendern.
  \item \textbf{Wiese/Blaetter}: Gruen-Saettigung erhoehen, Detail-Schaerfe verbessern.
  \item \textbf{Hintergrund-Unschaerfe}: Im Portraet-Modus erkennt die KI die Kontur des Hauptmotivs; der Hintergrund wird unscharf, das Motiv scharf.
\end{itemize}

\subsubsection*{Anwendungen des Deep Learning:}

\begin{itemize}
  \item \textbf{KI-Szenenerkennung}: Die Kamera erkennt automatisch Szenen wie „Essen", „Haustiere", „Nacht" oder „Dokument" und optimiert die Parameter entsprechend. Der Essen-Modus verstaerkt z. B. warme Toene und Saettigung, damit das Essen appetitlicher wirkt.
  \item \textbf{KI-Portraet-Freistellung}: Smartphone-Objektive koennen physisch keine so geringe Schaerfentiefe wie eine Spiegelreflexkamera erreichen; per Deep Learning wird die Motivkante modelliert und der Hintergrund unscharf gemacht. Neuere Algorithmen koennen sogar Haare und Glasrnder korrekt verarbeiten.
  \item \textbf{KI-Entfernung / KI-Bildweiterung}: Waehle Passanten oder Stoerobjekte im Bild aus; die KI „fuellt" den verdeckten Bereich anhand des Kontexts auf. Das ist ein Einsatz generativer KI in der Bildgebung.
\end{itemize}

\subsubsection*{Computational Photography vs. optische Fotografie:}

\begin{itemize}
  \item \textbf{Vorteile der Computational Photography}: Auch bei kompakten Abmessungen gute Bildqualitaet; auch in extremen Situationen (Schwachlicht, Gegenlicht) brauchbare Ergebnisse; normale Nutzer brauchen kein Fotowissen, um „gut aussehende" Fotos zu machen.
  \item \textbf{Nachteile der Computational Photography}: Manchmal „ueberarbeitet" – ueberschaerft, Farben verfaelscht, unnatuerlicher Himmel; bei Bewegung des Motivs entstehen bei Multi-Frame-Compositing Geisterbilder; Details sind „berechnet" statt real erfasst.
  \item \textbf{Vorteile der optischen Fotografie}: Grosser Sensor + lichtstarke Linse erfassen das Licht authentisch, natuerliche Bildqualitaet, grosser Spielraum fuer Nachbearbeitung; keine Algorithmus-„Halluzinationen".
  \item \textbf{Trend}: Moderne Flaggschiff-Smartphones kombinieren „Optik + Computation" – grosser Sensor fuer reale Daten, Algorithmen fuer die optimierte Darstellung. Beide Ansaetze sind nicht gegensaetzlich.
\end{itemize}

\booksubsection{VII. Objektivtypen und Brennweiten}

Moderne Flaggschiff-Smartphones haben meist drei bis vier Objektive, jedes fuer eine andere Brennweite. Die Kenntnis der Brennweiten hilft dir, in jeder Situation das richtige Objektiv zu waehlen.

\subsubsection*{Haeufige Objektive und aequivalente Brennweiten:}

\begin{itemize}
  \item \textbf{Ultraweitwinkel (13-24 mm aequivalent)}: Groesstes Sichtfeld, ideal fuer Landschaft, Architektur und Innenraeume. Erfasst grosse Szenen, aber die Raender verzerren (gerade Linien werden gekruemmt); bei Portraets werden Gesichter laenglich. Gut fuer „grossraeumige Szenen", ungeeignet fuer „Nahaufnahmen".
  \item \textbf{Hauptkamera (24-28 mm aequivalent)}: Die am haeufigsten genutzte Linse, beste Bildqualitaet und groesste Blende. Kompromiss aus Sichtfeld und Qualitaet, fuer Portraet- und Landschaftsaufnahmen. Bei den meisten Smartphones ist „1x" die Hauptkamera.
  \item \textbf{Tele (70-200 mm aequivalent)}: Fuer entfernte Motive und Portraet-Detailaufnahmen. Tele komprimiert die Raumwirkung (entfernte Berge und nahe Personen wirken naher beieinander); bei Portraets entstehen natuerlichere Gesichtsproportionen und Hintergrundunschaeerfe.
  \item \textbf{Periskop-Tele (ab 200 mm aequivalent)}: Durch gefaltete Optik wird in der Dicke eines Smartphones eine laengere Brennweite realisiert. Kann auf grosse Entfernung Details scharf abbilden – ideal fuer Konzerte, Sportevents und Tierfotografie. Nur in Flaggschiff-Geraeten.
\end{itemize}

\subsubsection*{Kompromiss zwischen Sensorgroesse und Pixelgroesse:}

Der Sensor ist die „Basis“ und bestimmt, wie viel Licht erfasst wird. In der Branche heisst es „ein groesserer Sensor schlaegt alles” – bei gleicher Pixelzahl bedeutet ein groesserer Sensor auch groessere Einzelpixel, mehr Licht und bessere Bildqualitaet.

\begin{itemize}
  \item \textbf{Grosser Sensor}: Viel Lichteinfall, gute Nachtbildqualitaet, breiter Dynamikumfang. Nachteil: Das Smartphone wird dicker (Kameramodul ragt heraus), hoeherer Preis.
  \item \textbf{Hohe Pixelzahl, kleiner Sensor}: Bei guter Tagesbeleuchtung liefern viele Pixel schaerfere Fotos und ermoeglichen spaeteres Cropping. Bei Schwachlicht erhaelt jeder Pixel zu wenig Licht – mehr Rauschen.
  \item \textbf{Massgebliche Flaggschiff-Wahl}: Die Hauptkamera liegt meist bei 1/1,3 bis 1 Zoll (einige Modelle verwenden bereits 1-Zoll-Sensoren) mit ca. 50 Megapixel – das ist der Kompromiss aus Bildqualitaet und Bautiefe.
\end{itemize}

\subsubsection*{Pixel-Binning-Technologie:}

Um hohe Pixelzahl bei Tag und grosse Pixel bei Nacht zu vereinen, verwenden Hersteller „Pixel-Binning“ – mehrere kleine Pixel werden zu einem grossen Ausgabepixel zusammengefasst.

\begin{itemize}
  \item \textbf{Quad-Bayer (4-in-1-Pixel)}: 50-MP-Sensor mit 12,5-MP-Ausgabe. Jedes Ausgabepixel wird aus 4 benachbarten Pixeln synthetisiert; Lichteinfall vervierfacht sich, sauberere Bildqualitaet.
  \item \textbf{Nonacell (9-in-1-Pixel)}: Fuer Sensoren ueber 100 MP; Ausgabe ca. 11-12 MP. Dasselbe Prinzip, noch groesserer Lichteinfall.
  \item \textbf{: Praktisch, aber mit niedrigem Wirkungsgrad und starker Erwärmung. Geeignet zum beiläufigen Aufladen im Büro, weniger geeignet, wenn es eilt, aus dem Haus zu kommen.}: Im Standardaufnahmemodus erhaeltst du das gebinnte niedrigapixelige Foto – gute Qualitaet, kleine Datei. Fuer maximale Details (z. B. Dokumente, Fernaufnahmen) kann manuell in den „Hochpixel-Modus" gewechselt werden, um die volle Aufloesung zu erhalten.
\end{itemize}

\booksubsection{VIII. Spezielle Aufnahmemodi}

Zusaetzlich zur Alltagsfotografie gibt es im Smartphone einige spezielle Modi, die, richtig eingesetzt, Aufnahmen ermoeglichen, die sonst nicht moeglich waeren.

\subsubsection*{Panorama-Modus:}

\begin{itemize}
  \item Beim Aufnehmen das Smartphone langsam und gleichmaessig in Pfeilrichtung quer bewegen; nicht zu schnell – sonst entstehen Nahtfehler.
  \item Das Smartphone vertikal stabil halten und nur horizontal verschieben. Bei Auf- und Abbewegung wird der Horizont wellig.
  \item Bei sich bewegenden Personen oder Fahrzeugen vorsichtig einsetzen – beim Zusammenfuegen koennen „halbe Koerper" oder „Geisterbilder" entstehen.
  \item Panorama kann auch vertikal aufgenommen werden – vertikal fotografieren, horizontal zusammenfuegen; ideal fuer hohe Gebaeude oder Wasserfaelle.
\end{itemize}

\subsubsection*{Slow-Motion und Zeitraffer:}

\begin{itemize}
  \item \textbf{Slow-Motion}: Aufnahme mit 120 fps oder 240 fps, bei der Wiedergabe als Zeitlupe. Ideal fuer Wassertropfen, Sportmomente oder springende Tiere. Ausreichend Licht erforderlich – bei Slow-Motion ist die Belichtungszeit kurz, in dunkler Umgebung ist die Qualitaet schlecht.
  \item \textbf{Zeitraffer}: In Intervallen aufgenommene Fotos werden zu einem Zeitraffer-Video zusammengefuegt. Ideal fuer Verkehr, ziehende Wolken, Sonnenuntergaenge oder bluehende Blueten. Aufnahmedauer von Minuten bis Stunden; zwingend Stativ verwenden, da aus der Hand unscharf.
  \item \textbf{Tipp}: Beim Zeitraffer von Wolken das Smartphone 15-30 Minuten aufstellen; das Ergebnis sind etwa 10 Sekunden Zeitraffer-Video mit beeindruckender Wirkung.
\end{itemize}

\subsubsection*{Sternenhimmel-/Astronomie-Modus:}

\begin{itemize}
  \item Abseits der Lichtverschmutzung der Stadt aufhalten; ein Vorort oder Gebirge ohne Strassenlaterne ist notwendig. In der Stadt ist die Milchstrasse nicht erfassbar.
  \item Zwingend Stativ verwenden – aus der Hand wird es unscharf. Das Smartphone muss einige Sekunden bis zu mehreren Zehger Sekunden ruhig stehen.
  \item Klarer Himmel ohne Wolken, der Mond sollte nicht zu gross sein (Vollmond ueberstrahlt die Sterne). Die dunkelsten Naechte um Monatsanfang oder -ende sind am besten.
  \item Der „Nachtmodus" des iPhone erkennt den Sternenhimmel automatisch und verlaengert die Belichtungszeit (bis zu 30 Sekunden); der „Sternenhimmel-Modus" oder „Astronomie-Modus" bei Android-Flaggschiffen funktioniert aehnlich. Nach dem Ausloesen das Smartphone nicht bewegen.
\end{itemize}

\subsubsection*{Makro-Modus:}

\begin{itemize}
  \item Fuer Blueten, Insekten, Schmuck oder Textdetailaufnahmen. Fokus bis auf 1-2 cm an das Motiv heran.
  \item Einige Smartphones haben eine dedizierte Makrolinse (meist 2 oder 5 MP mit kleinem Sensor) mit mittelmassiger Qualitaet; Flaggschiffe nutzen die Ultraweitwinkellinse auch fuer Makro, mit besserer Qualitaet.
  \item Beim Nah-Aufnehmen ist die Linse sehr nah am Motiv und verdeckt Licht; bei Bedarf mit einer Schreibtischlampe oder Taschenlampe Zusatzlicht geben.
  \item Bei extrem kurzer Fokusdistanz ist die Schaerfentiefe sehr gering; das Motiv kann bei wenigen Millimetern Bewegung unscharf werden. Ruehige Hand und unbewegtes Motiv sind wichtig.
\end{itemize}

\begin{tcbinfo}{Praktische Tipps fuer Standfotografie}
\begin{itemize}
  \item Alltag: HEIF/JPEG ist ausreichend – kleine Dateien, brauchbare Qualitaet.
  \item Wichtige Motive (Landschaft, Portraet, Nachbearbeitung): RAW aktivieren, um maximalen Spielraum fuer die Bearbeitung zu behalten.
  \item Richtiges Objektiv: Landschaft = Ultraweitwinkel, Portraet = Tele, Alltag = Hauptkamera.
  \item Nacht und Sternenhimmel: Smartphone zwingend fixieren, Nacht- oder Astronomie-Modus aktivieren.
  \item Grosssensorgeraete (z. B. mit 1-Zoll-Sensor) haben ein hoeheres Qualitaetspotenzial, aber die Bedienung bleibt entscheidend.
\end{itemize}

\end{tcbinfo}

\begin{tcbpractice}{Probier es aus: Lerne deine Kamera kennen}
\begin{enumerate}
  \item Pruefe, ob dein Smartphone RAW-Aufnahmen unterstuetzt (Einstellungen Kamera Formate), aktiviere es, mache eine Aufnahme und vergleiche die Dateigroessen von JPEG und RAW.
  \item Wechsle zwischen Ultraweitwinkel, Hauptkamera und Tele, fotografiere dasselbe ferne Motiv und vergleiche Sichtfeld und Qualitaet.
  \item Finde eine Szene mit Nacht- oder Sternenhimmel-Modus, fixiere das Smartphone und mache eine Nachtaufnahme.
  \item Versuche, im Panorama-Modus eine breite Szene quer aufzunehmen; das Smartphone gleichmaessig bewegen.
\end{enumerate}

\end{tcbpractice}
